\section{Cơ sở lý thuyết}

Chương này gom các khái niệm cần dùng khi đọc PCA, KPCA và các phương pháp học manifold. Phần lớn là đại số ma trận quen thuộc; ta chỉ nhắc lại đủ để các chứng minh ở Chương~3 mạch lạc.

\subsection{Ma trận dữ liệu và hiệp phương sai}

Với mẫu $x_1,\ldots,x_m \in \mathcal{X}$, ma trận $\mathbf{X}$ như Chương~1. Khi nói dữ liệu \emph{đã căn giữa theo cột}, ta có $\sum_{i=1}^m \mathbf{x}_i = \mathbf{0}$.

\begin{definition}[Hiệp phương sai mẫu]
Với dữ liệu đã căn giữa,
\[
\mathbf{C} = \frac{1}{m}\,\mathbf{X}\mathbf{X}^\top \in \R^{N \times N}.
\]
Phần tử $\mathbf{C}_{ij}$ đo mức đồng biến tuyến tính giữa hai đặc trưng sau khi trừ trung bình.
\end{definition}

Nếu chưa căn giữa, dùng $\bar{\mathbf{x}} = \frac{1}{m}\sum_i \mathbf{x}_i$ và
$\mathbf{C} = \frac{1}{m}\sum_i (\mathbf{x}_i - \bar{\mathbf{x}})(\mathbf{x}_i - \bar{\mathbf{x}})^\top$.

\subsection{Phân tích giá trị kỳ dị}

\begin{definition}[SVD]
Với $\mathbf{X} \in \R^{N \times m}$, tồn tại phân tích
\[
\mathbf{X} = \mathbf{U}\boldsymbol{\Sigma}\mathbf{V}^\top,
\]
$\mathbf{U},\mathbf{V}$ có cột trực chuẩn, $\boldsymbol{\Sigma}$ đường chéo các giá trị kỳ dị không âm $\sigma_1 \geq \sigma_2 \geq \cdots$.
\end{definition}

Đặt $\boldsymbol{\Lambda} = \boldsymbol{\Sigma}^2$. Khi đó
\[
\mathbf{C} = \frac{1}{m}\mathbf{U}\boldsymbol{\Lambda}\mathbf{U}^\top,
\qquad
\mathbf{K} = \mathbf{X}^\top\mathbf{X} = \mathbf{V}\boldsymbol{\Lambda}\mathbf{V}^\top,
\]
với $\mathbf{K}$ là ma trận Gram (kernel tuyến tính trên các cột $\mathbf{x}_i$). Cùng một bộ giá trị kỳ dị $\lambda_i = \sigma_i^2/m$ xuất hiện ở $\mathbf{C}$ và $\mathbf{K}$, chỉ khác hệ vector kỳ dị trái/phải---đây là cầu nối giữa PCA và KPCA.

\subsection{Chiếu trực giao và chuẩn Frobenius}

\begin{definition}[Ma trận chiếu bậc $k$]
$P_k$ gồm các ma trận $\mathbf{P} = \mathbf{U}\mathbf{U}^\top$ với $\mathbf{U} \in \R^{N \times k}$ có cột trực chuẩn, $\rank(\mathbf{P}) = k$. Khi đó $\mathbf{P}^2 = \mathbf{P}$, $\mathbf{P}^\top = \mathbf{P}$.
\end{definition}

Chuẩn Frobenius: $\|\mathbf{M}\|_F^2 = \Tr(\mathbf{M}^\top\mathbf{M})$. Tính hoán vị của trace ($\Tr(\mathbf{A}\mathbf{B}\mathbf{C}) = \Tr(\mathbf{B}\mathbf{C}\mathbf{A})$) dùng nhiều trong chứng minh PCA.

\subsection{Phương sai theo một hướng}

Với vector đơn vị $\mathbf{u}$, lượng $\mathbf{u}^\top\mathbf{C}\mathbf{u}$ là phương sai của chiếu một chiều lên $\mathbf{u}$. Bài toán $\max_{\|\mathbf{u}\|=1} \mathbf{u}^\top\mathbf{C}\mathbf{u}$ đạt tại vector kỳ dị của $\mathbf{C}$ ứng với giá trị kỳ dị lớn nhất (thương số Rayleigh). Các hướng tiếp theo tối đa hóa phương sai trong không gian vuông góc với các hướng đã chọn.

\begin{example}
Hai đặc trưng trên mặt phẳng tương quan mạnh (ví dụ cùng đo kích thước giày nhưng đơn vị khác nhau). Đám điểm kéo dài theo một đường chéo; chiếu lên đường đó giữ phần lớn biến thiên, chiếu vuông góc gần như mất thông tin---đó là động lực hình học của thành phần chính thứ nhất (Hình~\ref{fig:pca-shoe}).
\end{example}

\begin{figure}[H]
\centering
\begin{tikzpicture}[scale=1.0]
  \draw[->, thick] (-4.5,0) -- (4.5,0) node[below] {$x^{(1)}$};
  \draw[->, thick] (0,-4) -- (0,4.2) node[left] {$x^{(2)}$};
  \foreach \x/\y in {-3.5/-2.8, -2/-1.6, 0/0, 2/1.6, 3.5/2.8, -3/-2.4, 2.6/2.1} {
    \fill[blue!70] (\x,\y) circle (2.5pt);
  }
  \draw[red, very thick] (-4,-3.2) -- (4,3.2);
  \node[red, font=\small] at (3.2,3.6) {PC1};
\end{tikzpicture}
\caption{Dữ liệu hai chiều (điểm xanh) và hướng thành phần chính thứ nhất (đường đỏ).}
\label{fig:pca-shoe}
\end{figure}

\subsection{Kernel và ma trận Gram}

\begin{definition}[Kernel xác định dương bán]
Hàm $K : \mathcal{X} \times \mathcal{X} \to \R$ là kernel PDS nếu tồn tại $\mathcal{H}$ và $\Phi : \mathcal{X} \to \mathcal{H}$ sao cho $K(x,x') = \langle \Phi(x), \Phi(x') \rangle$.
\end{definition}

Trên $m$ điểm, ma trận $\mathbf{K}$ với $K_{ij} = K(x_i,x_j)$ luôn xác định dương bán. KPCA chỉ cần $\mathbf{K}$, không cần biết $\Phi$ tường minh.

\subsection{Đồ thị láng giềng và Laplacian}

Các thuật toán manifold thường:
\begin{enumerate}
  \item Chọn $t$ láng giềng gần nhất cho mỗi điểm (theo $\|\mathbf{x}_i - \mathbf{x}_j\|$).
  \item Gán trọng số $W_{ij}$ trên cạnh (Gaussian, hệ số tái tạo,\ldots).
  \item Dùng ma trận Laplacian $\mathbf{L} = \mathbf{D} - \mathbf{W}$, với $D_{ii} = \sum_j W_{ij}$.
\end{enumerate}

Ma trận căn giữa (centering) dùng trong Isomap:
\[
\mathbf{H} = \mathbf{I}_m - \frac{1}{m}\mathbf{1}\mathbf{1}^\top,
\]
với $\mathbf{1}$ là vector toàn $1$. Nhân $\mathbf{H}$ hai bên một ma trận khoảng cách giúp chuyển sang dạng tương tự tích vô hướng trong không gian đã trừ trung bình.

\FloatBarrier
